\medskip
\noindent\textbf{Abstract.}\par\smallskip
This report reviews mathematical models and numerical methods used to simulate
blood flow, with emphasis on large-vessel and stenotic settings. The review
starts from continuum balance laws and incompressible Navier--Stokes models,
then compares three-dimensional CFD, fluid--structure interaction, axisymmetric
and two-dimensional reductions, distributed one-dimensional area--flow models,
lumped zero-dimensional networks, geometrical multiscale coupling, and emerging
learned surrogate approaches. The review is organized around source-to-source
contrasts: which equations and state variables are retained, which physical
effects are closed, which data are required, which numerical evidence is
reported, and which observables are native or reconstructed. Across those
tiers, the report separates the modeling choices that determine interpretation:
rheology, wall mechanics, geometry, inlet and outlet data, numerical
discretization, verification, validation, and observable definitions.

The synthesis compares model classes across fidelity, cost, retained state,
observables, identifiability, validation status, reproducibility, and
uncertainty. Its central conclusion is that no model tier is interpretable in
isolation: equations, closure assumptions, data, numerical evidence, and
observation maps must be matched to the question being asked. An idealized
stenosis computation is then used as a worked example of those review
principles. The case study uses an $R_{\max}$-normalized Canic-derived
one-dimensional area--flow model with solver variables $a=R^2$ and
$q=Q_{\mathrm{phys}}/\pi$. It records the model contract, closure choices,
finite-volume realization, manufactured-solution verification record,
geometry-rest equilibrium test, production diagnostics, and a
plane--tetrahedron observation operator for comparing one-dimensional output
with available resolved three-dimensional velocity data. The MMS record gives
bounded implementation-verification evidence; the zero-inlet rest audit shows
that the scheme is not exactly well balanced; and the timestamp-matched
1D--3D section comparison reports descriptive velocity discrepancies under the
declared observation operator. Because the resolved wall, boundary, material,
geometry-state, and transient-history records are not fully matched, those
differences are not validation errors. Radial-profile output is excluded
pending operator-test and reducer reconciliation. The case study illustrates
how model reduction, closure choices, equilibrium preservation, and matched
observation operators affect the meaning of computed blood-flow quantities.
